% =====================================================================
%  The MOND acceleration scale as the cosmic dynamical acceleration
%  Carl Zimmerman, 2026.  Self-contained REVTeX 4-2 source for Zenodo/arXiv.
%  Compile:  pdflatex Zimmerman_Scaling_MOND_2026.tex   (twice, for refs)
%  All numerical results are reproducible from the public repository cited
%  in the Data Availability statement.
% =====================================================================
\documentclass[aps,prd,preprint,nofootinbib,floatfix]{revtex4-2}
\usepackage{amsmath,amssymb}
\usepackage{url}

\begin{document}

\title{The MOND acceleration scale as the cosmic dynamical acceleration:\\
a scaling--MOND cosmology and its bridges to the dark sector}

\author{Carl Zimmerman}
\affiliation{Independent Researcher}
\date{\today}

\begin{abstract}
We examine the proposal that the Milgromian acceleration scale is the cosmic
dynamical acceleration,
$a_0=\tfrac{c}{2}\sqrt{G\rho_c}=cH(z)/Z$ with $Z=2\sqrt{8\pi/3}\simeq5.789$.
The evolving-scale hypothesis itself is Milgrom's (2014); our contribution is an
explicit covariant realization of it and a proof of its CMB-safety. Algebraically
$a_0=cH/Z$ re-expresses the long-known coincidence $a_0\approx cH_0$ through the
Friedmann relation; its coefficient-free consequence is that the scale {\it evolves},
$a_0(z)=a_0(0)\,E(z)$. Fitting $a_0\propto E(z)^p$ to the recent compilation (SPARC;
V\u{a}r\u{a}\c{s}teanu 2025; MUSE--DARK III 2026) gives $p=0.80\pm0.17$; constant $a_0$
is excluded at $5\sigma$ in the nominal fit, but a jackknife and an inter-method
systematic (the $1.7\sigma$ discrepancy between the two near-coincident local anchors)
reduce this to $\simeq2\sigma$, so we treat the trend as a suggestive hint, not a
detection. Working from the Aether--Scalar--Tensor (AeST) action we promote the scale
to the aether expansion, $a_0=c\theta/3Z$, and show by order-counting---reinforced by
the identity $\delta q^{00}=0$ that survives perturbing the aether itself---that $a_0$
is absent from the {\it linear} cosmological perturbations; the running therefore
leaves the linear CMB and $P(k)$ exactly invariant, which we confirm by direct
integration of the Einstein--Boltzmann system ($r_s=144.3$~Mpc, $\ell_A=301.7$). We
are explicit about the boundaries: the $O(1)$ coefficient $Z$ is a posit; the
Standard-Model constants do not follow; and the apparent-$a_0$ evolution is itself
reproduced by $\Lambda$CDM hydrodynamic simulations, so the amplitude alone does not
discriminate. The one prediction that {\it does}---a redshift-weakening external field
effect, $\eta=g_{\rm ext}/a_0(z)\propto1/E(z)$, distinct from both $\Lambda$CDM and
constant-$a_0$ MOND---is forecast to require a dedicated $z\gtrsim4$ campaign.
\end{abstract}

\maketitle

% ---------------------------------------------------------------------
\section{Introduction}
The Milgromian scale $a_0\simeq1.2\times10^{-10}\,{\rm m\,s^{-2}}$, below which
galaxy dynamics depart systematically from the Newtonian expectation, is
numerically close to the cosmic acceleration $cH_0$~\cite{Milgrom1983,Famaey2012}.
This proximity, $a_0\approx cH_0/(2\pi)$, has been noted for forty years and is
usually treated as either a coincidence or a clue that $a_0$ is set by
cosmology~\cite{Milgrom2011}. Here we take the second view in a specific form,
write the scale as a cosmic dynamical acceleration, derive its one falsifiable
consequence, confront that with data, and map---honestly---how far the picture
can be pushed toward a relativistic theory of the dark sector. We are equally
explicit (Sec.~\ref{sec:scope}) about what the construction does {\it not} do.

% ---------------------------------------------------------------------
\section{The premise and its coefficient}
We write the acceleration scale as half the surface gravity of the cosmic
free-fall density,
\begin{equation}
a_0=\frac{c}{2}\sqrt{G\rho_c}=\frac{c^2}{2R},\qquad
R=\sqrt{\tfrac{8\pi}{3}}\,\frac{c}{H}=c\,t_{\rm ff},
\label{eq:premise}
\end{equation}
where $t_{\rm ff}=(G\rho_c)^{-1/2}$ is the dynamical time of the critical-density
medium. Using the Friedmann relation $H^2=\tfrac{8\pi}{3}G\rho_c$,
Eq.~(\ref{eq:premise}) is equivalent to
\begin{equation}
a_0=\frac{cH}{Z},\qquad
Z=2\sqrt{\tfrac{8\pi}{3}}=\sqrt{\tfrac{32\pi}{3}}\simeq5.789 .
\label{eq:Z}
\end{equation}
The factor $\sqrt{8\pi/3}$ is the Friedmann free-fall-to-Hubble ratio and is not
free; the remaining factor of two---the Schwarzschild reading $c^2/2R$---is the
only undetermined element. The density/surface-gravity {\it form} of
Eq.~(\ref{eq:premise}) appears to be non-standard (the literature uses $cH_0$ or
$\sqrt{\Lambda}$, not $\sqrt{G\rho}$), but it is a re-expression of a known
coincidence rather than new dynamics, and the coefficient is not unique: the
de~Sitter-horizon readings of Milgrom ($Z=2\pi$) and Verlinde
($Z\simeq6$)~\cite{Verlinde2017} fit the same data equally well. We return to the
status of $Z$ in Sec.~\ref{sec:coeff}.

% ---------------------------------------------------------------------
\section{The falsifiable prediction and the data}
\label{sec:data}
Because $a_0$ tracks a cosmic density, it evolves---a possibility Milgrom himself
raised~\cite{Milgrom2014}; and because $Z$ cancels in the ratio, the prediction is
coefficient-free and independent of the Hubble tension:
\begin{equation}
\frac{a_0(z)}{a_0(0)}=E(z)=\sqrt{\Omega_m(1+z)^3+\Omega_\Lambda}.
\label{eq:Ez}
\end{equation}
Testing such evolution through the Tully--Fisher zero point was proposed early
as a modified-gravity probe by Gnedin~\cite{Gnedin2008}.
We fit $a_0(z)=A\,E(z)^p$ to the compilation in Table~\ref{tab:data},
marginalizing the normalization $A$. The exponent is $p=0.80\pm0.17$ ($1\sigma$);
the model comparison gives $\chi^2=27$ for constant $a_0$, $3.8$ for $E(z)$, and
$27.5$ for the matter-only law $(1+z)^{3/2}$. Thus in the nominal fit constant $a_0$ and the
matter-only branch are each excluded at $\simeq5\sigma$, while $p=1$ is consistent
at $1.1\sigma$. This significance is fragile, and we state so plainly. A jackknife
shows the rejection of constant $a_0$ rests almost entirely on the single
$z\simeq0.9$ point---dropping it leaves only $1.2\sigma$---and the $1.7\sigma$
discrepancy between the two near-coincident local anchors ($1.20$ versus $1.69$ at
$z\simeq0$, which cannot reflect genuine evolution over $\Delta z=0.05$) evidences an
inter-method systematic $\sigma_{\rm sys}\simeq0.28$ (same units); folding
$\sigma_{\rm sys}$ into all points reduces the rejection of constant $a_0$ to
$\simeq2\sigma$. We therefore regard the trend as a suggestive hint, not a detection,
dominated by the local-anchor systematic rather than the redshift law. A single clean
deep-MOND measurement at $z=3$ with $3\%$ precision would sharpen the exponent to
$p\simeq0.99\pm0.04$.

We emphasize two honest caveats. First, an evolving radial-acceleration relation is
{\it also} expected in $\Lambda$CDM: hydrodynamic simulations (Magneticum,
Tian~\textit{et~al.}~\cite{Tian2022}) find the apparent best-fit $a_0$ rising by
$\simeq3$ from $z=0$ to $z=2.3$---matching $E(2.3)=3.46$---with no fundamental
evolution. The data thus favor ``evolving over constant,'' not the present framework
over $\Lambda$CDM. Second, the dispersion-based dynamical masses of de
Graaff~\textit{et~al.}~\cite{deGraaff2024} at $z\simeq6$ reach $M_{\rm dyn}/M_\star\sim40$;
these are {\it not} accounted for by the evolving scale, whose deep-MOND boost
$\sqrt{a_0(z)/g_{\rm bar}}$ reaches only $\sim3$ for such compact systems (which lie
near the Newtonian regime at their epoch, $g_{\rm bar}\gtrsim a_0(z)$). They are a
puzzle for MOND, evolving or not, and we do not cite them as support.

\begin{table}[t]
\caption{The $a_0(z)$ compilation used in the fit. $a_0$ in
$10^{-10}\,{\rm m\,s^{-2}}$. Provenance is tabulated in the repository
(Data Availability).}
\label{tab:data}
\begin{ruledtabular}
\begin{tabular}{ccl}
$z$ & $a_0$ & source\\
\hline
$0.00$ & $1.20\pm0.26$ & SPARC---McGaugh, Lelli \& Schombert 2016~\cite{McGaugh2016}\\
$0.05$ & $1.69\pm0.13$ & V\u{a}r\u{a}\c{s}teanu \textit{et al.} 2025~\cite{Varasteanu2025}\\
$0.90$ & $2.38\pm0.11$ & MUSE--DARK III 2026~\cite{MUSEDARK2026}\\
\end{tabular}
\end{ruledtabular}
\end{table}

% ---------------------------------------------------------------------
\section{The over-constrained web}
The content of Eq.~(\ref{eq:Z}) is a single dimensionless invariant,
\begin{equation}
\frac{a_0}{cH(z)}=\frac{1}{Z}=0.173 \quad\text{(every epoch)},
\end{equation}
projected by dimensional analysis into different units---an acceleration (the
radial-acceleration relation), a length $\ell=c^2/a_0=Z R_H$, a surface density
$\Sigma=a_0/G$, a velocity (the baryonic Tully--Fisher zero point), a temperature
$T_{\rm dS}/Z$, and a time (the cosmic free-fall time). From the one premise about
a dozen relations follow as forced consequences (e.g. $H_0=Za_0/c$; the de~Sitter
floor $a_0(0)\sqrt{\Omega_\Lambda}$; the Tully--Fisher and Faber--Jackson
zero-points $\propto1/E(z)$). The structure is {\it over-constrained}: the same
$a_0$ must simultaneously fit the SPARC scale, the $\Lambda$ floor
$a_0(0)\sqrt{\Omega_\Lambda}$, and the intermediate-redshift points
(V\u{a}r\u{a}\c{s}teanu, MUSE--DARK)---several determinations across cosmic time keyed
to one number (the relation $H_0=Za_0/c$ is an identity, not an independent check). This motivates a sharp filter: a
genuine relation off Eq.~(\ref{eq:Z}) is $z$-invariant ($a_0/cH$ stays $1/Z$),
whereas a coincidence between cosmic numbers is not. The same test that retains
these edges {\it rejects}, for example, the present-epoch identity
$2\sin^2\theta_W\simeq\Omega_m/\Omega_\Lambda$, whose ratio drifts as
$(1+z)^{-3}$---a why-now coincidence, not a law.

% ---------------------------------------------------------------------
\section{Relativistic completion and CMB-safety}
\label{sec:bridge1}
A covariant home is the Aether--Scalar--Tensor (AeST) theory of Skordis \&
Zlosnik~\cite{Skordis2021}, the one relativistic MOND theory that reproduces the
CMB and matter power spectra. Its action is
\begin{multline}
S=\!\int\! d^4x\,\frac{\sqrt{-g}}{16\pi\tilde G}\Big[R-\tfrac{K_B}{2}F^{\mu\nu}F_{\mu\nu}
+2(2-K_B)J^\mu\nabla_\mu\phi\\
-(2-K_B)\mathcal{Y}-\mathcal{F}(\mathcal{Y},\mathcal{Q})
-\lambda(A^\mu A_\mu+1)\Big]+S_m[g],
\label{eq:aest}
\end{multline}
with a unit-timelike vector ($A^\mu A_\mu=-1$), the spatial and temporal kinetic
scalars $\mathcal{Y}=q^{\mu\nu}\nabla_\mu\phi\nabla_\nu\phi$ and
$\mathcal{Q}=A^\mu\nabla_\mu\phi$ ($q_{\mu\nu}=g_{\mu\nu}+A_\mu A_\nu$), and a free
function $\mathcal{F}$. The MOND scale enters the {\it spatial} sector,
$\mathcal{F}\to[2\lambda_s/(1+\lambda_s)a_0]\,\mathcal{Y}^{3/2}$ as
$\nabla\phi\to0$, whereas the CMB-fitting, dark-matter-mimicking mode is carried
by the {\it temporal} sector,
$\mathcal{K}(\mathcal{Q})=-2\Lambda+\mathcal{K}_2(\mathcal{Q}-\mathcal{Q}_0)^2$,
which gives $\bar\rho\propto a^{-3}$ and contains no $a_0$.

We promote the constant $a_0$ to $a_0(z)=c\,\theta/(3Z)$, where
$\theta=\nabla_\mu A^\mu$ is the aether expansion ($=3H$ on a
Friedmann--Robertson--Walker background); this reproduces Eq.~(\ref{eq:Z}) and is
a minimal, covariant modification because $\theta$ is already present in
Eq.~(\ref{eq:aest}). The key result is an order-counting statement. On the
homogeneous background $\bar\phi=\bar\phi(t)$ is purely temporal, and because
$q^{\mu\nu}$ projects orthogonal to $A^\mu$ one has $\bar{\mathcal{Y}}=0$ and
$\delta\mathcal{Y}=0$ at linear order. The one term the standard order-counting
omits---$\delta q^{\mu\nu}\nabla_\mu\bar\phi\,\nabla_\nu\bar\phi$, surviving only as
$\delta q^{00}(\dot{\bar\phi})^2$---also vanishes identically, because the
unit-timelike constraint forces $\delta A^0=-\Psi$ and hence
$\delta q^{00}=\delta g^{00}+2A^0\delta A^0=2\Psi-2\Psi=0$. Thus robustly
\begin{equation}
\mathcal{Y}=O(\delta\phi^2),\qquad
\mathcal{Y}^{3/2}=O(\delta\phi^3).
\end{equation}
The $a_0$-bearing term is therefore third order in perturbations: it makes no
contribution to the second-order action and hence none to the {\it linear}
equations of motion, which are governed by the analytic term
$-(2-K_B)\mathcal{Y}$ and by $\mathcal{K}(\mathcal{Q})$, both $a_0$-independent.
Consequently
\begin{equation}
a_0\to a_0(z)\ \Rightarrow\ \{C_\ell,\,P(k)\}_{\rm linear}\ \text{exactly invariant.}
\end{equation}
We verified this with a direct integration of the linear Einstein--Boltzmann
system: the running-$a_0$ effect on the linear spectra is numerically zero, and
the background acoustic physics is recovered correctly ($r_s=144.3$~Mpc,
$\ell_A=301.7$, against the Planck values $144.4$~Mpc and $\simeq301$). The
running of $a_0$ thus enters only at nonlinear/quasi-static order. A
Planck-precision $C_\ell$ amplitude calculation with the full vector sector is
left to a dedicated Boltzmann code; it would reproduce the known AeST fit and does
not alter the linear invariance just established. The $\mathcal{Y}^{3/2}$
non-analyticity makes the nonlinear matching delicate; the leading order-counting
is robust.

% ---------------------------------------------------------------------
\section{Dark-sector unification, and the coefficient}
\label{sec:bridge3}
In Eq.~(\ref{eq:aest}) the single temporal function $\mathcal{K}(\mathcal{Q})$
carries {\it both} dark sectors: the dust mode
($\mathcal{K}_2(\mathcal{Q}-\mathcal{Q}_0)^2\to\bar\rho\propto a^{-3}$, mimicking
cold dark matter) and the cosmological constant ($-2\Lambda$). The evolving scale
ties them: $a_0(z)$ floors at $a_0(0)\sqrt{\Omega_\Lambda}$, the $\Lambda$ value.
Dark matter and dark energy are, in this reading, one evolving scale at its two
ends rather than two substances. The claim is falsifiable through
$a_0$-cosmography: $H(z)=Z\,a_0(z)/c$ reconstructs the expansion from galaxy
dynamics. The present data give $H(0)\simeq71.5$ (within the $H_0$ band) and a
$\sim25\%$ excess at $z\simeq0.9$---the same anchor-driven tension seen in
Sec.~\ref{sec:data}. A clean $a_0$ at $z>2$ converts this into a measurement of
the dark-energy equation of state from the dark-matter sector.

\paragraph*{The coefficient.}\label{sec:coeff}
Horizon thermodynamics~\cite{Jacobson1995,Verlinde2017} robustly yields the
{\it evolution} $a_0\propto H$ and the order of magnitude, but does not pin the
$O(1)$ coefficient: distinct, individually defensible matchings give $Z=1$
(Unruh${=}$de~Sitter temperature), $Z=2\pi$ (Milgrom), $Z\simeq6$ (Verlinde), and
$Z=5.789$ (the present density form). The framework's ``Schwarzschild'' reading
is not self-consistent---the mass enclosed within $R$ is $8\pi/3$ times the
Schwarzschild mass for $R$, so $R$ is not a horizon---hence the factor of two is
a posit. The physically motivated readings cluster within $8.5\%$, inside the
$\sim20\%$ systematic on $a_0$, so the data cannot select among them. This costs
the framework nothing, because the falsifiable prediction, Eq.~(\ref{eq:Ez}), is
$Z$-independent.

% ---------------------------------------------------------------------
\section{Scope: what is not claimed}
\label{sec:scope}
The construction is a candidate description of the {\it dark sector}
(gravity, dark matter, dark energy), not a theory of everything, and we state the
boundary plainly because related notes have not. The Standard-Model constants do
{\it not} follow from Eq.~(\ref{eq:Z}). A search over $\sim3.4\times10^4$ simple
formulae built from $Z$ matches an arbitrary $O(100)$ target to $0.004\%$ roughly
$20\%$ of the time; $52$ of $64$ such ``derivations'' contain no $Z$ at all; and a
randomly chosen number reproduces $\alpha^{-1}$ and the lepton/baryon mass ratios
to the same precision as $32\pi/3$. Retrodictions of this kind carry $\approx0$
bits of evidence and form no part of the present framework. Likewise we do not
claim a first-principles value for $Z$ (Sec.~\ref{sec:coeff}) nor a
Planck-precision CMB fit (Sec.~\ref{sec:bridge1}). Stating these limits is what
keeps the surviving result---Eqs.~(\ref{eq:premise})--(\ref{eq:Ez}) and the
$5\sigma$ rejection of constant $a_0$---trustworthy.

% ---------------------------------------------------------------------
\section{The decisive test and conclusion}
Pinning the exponent $p$ by measuring $a_0$ at $z>2$ tests evolution against
constant $a_0$, but---because $\Lambda$CDM's {\it apparent} $a_0$ also rises as
$E(z)$ (Sec.~\ref{sec:data})---it does {\it not} by itself separate the present
framework from $\Lambda$CDM. The deep-MOND consequence cascade
($M_{\rm dyn}/M_\star\propto\sqrt{E}$, $\sigma\propto E^{1/4}$, the Tully--Fisher
zero point $\propto-\log E$) shares this degeneracy: every channel follows from the
radial-acceleration knee at $a_0(z)$, and $\Lambda$CDM, with apparent
$a_0\simeq a_0(0)E(z)$, returns the same powers.

The one observable that discriminates is the {\it external field effect}. With
$a_0\to a_0(z)$, a galaxy in a fixed environment $g_{\rm ext}$ has
$\eta=g_{\rm ext}/a_0(z)\propto1/E(z)$: the EFE {\it weakens} with redshift, so
high-$z$ galaxies in dense environments behave more like isolated deep-MOND systems.
This is distinct from $\Lambda$CDM (no host effect on the internal relation) {\it and}
from constant-$a_0$ MOND (constant EFE). Solving the MOND EFE numerically, the
distinctive signal---the change across redshift in the embedded-versus-isolated
mass-discrepancy offset---is $\simeq0.1$~dex, below the per-galaxy scatter
($\simeq0.25$--$0.4$~dex from kinematics, stellar and gas masses, and intrinsic
scatter). A $3\sigma$ detection requires $\sim600$--$1600$ {\it extended} (deep-MOND),
environment-classified, kinematically-resolved galaxies at $z\gtrsim4$---a dedicated
multi-cycle JWST/ALMA(/ELT) program, not a single pointing.

In summary: $a_0=\tfrac{c}{2}\sqrt{G\rho_c}=cH(z)/Z$ re-expresses Milgrom's
coincidence in a density form; its evolving consequence $a_0\propto E(z)$---an idea
due to Milgrom~\cite{Milgrom2014}---is realized here covariantly in AeST, is provably
CMB-safe at linear order, and is favored over constant $a_0$ by current data at
$\simeq2\sigma$. The coefficient remains a posit, the Standard-Model sector is
untouched, and the amplitude is $\Lambda$CDM-degenerate. The framework's one
distinctive, falsifiable signature is the redshift-weakening EFE, whose measurement we
have forecast.

% ---------------------------------------------------------------------
\begin{acknowledgments}
The analysis, numerical calculations, and drafting of this manuscript were carried
out with the assistance of Claude (Anthropic, 2026), used as a tool; the author is
responsible for all claims. This work makes use of the publicly available SPARC
database and of published high-redshift kinematic measurements as cited.
\end{acknowledgments}

\paragraph*{Data and code availability.}
All calculations, the $a_0(z)$ compilation with its per-point provenance, and the
linear Einstein--Boltzmann solver are reproducible from the public repository
\url{https://github.com/carlzimmerman/zimmerman-formula} (directory
\texttt{real\_research/}; scripts \texttt{a0\_decisive\_pipeline.py},
\texttt{REAL\_WEB.py}, \texttt{bridge1\_linear\_boltzmann.py},
\texttt{bridge1\_aest\_equations.md},
\texttt{bridge2\_coefficient\_thermodynamics.py},
\texttt{bridge3\_dark\_sector\_unification.py},
\texttt{false\_discovery\_rate.py}, \texttt{can\_another\_number\_do\_it.py}).

% ---------------------------------------------------------------------
\begin{thebibliography}{99}
\bibitem{Milgrom1983} M.~Milgrom, Astrophys.\ J.\ {\bf 270}, 365 (1983).
\bibitem{Milgrom2011} M.~Milgrom, ``MOND's acceleration scale as a fundamental
  quantity,'' arXiv:1110.2580; Mod.\ Phys.\ Lett.\ A {\bf 26}, 2677 (2011).
\bibitem{Milgrom2014} M.~Milgrom, ``Cosmological variation of the MOND constant:
  secular effects on galactic systems,'' Phys.\ Rev.\ D {\bf 91}, 044009 (2015);
  arXiv:1412.4344.
\bibitem{Famaey2012} B.~Famaey and S.~S.~McGaugh, Living Rev.\ Relativity
  {\bf 15}, 10 (2012).
\bibitem{McGaugh2016} S.~S.~McGaugh, F.~Lelli, and J.~M.~Schombert, Phys.\ Rev.\
  Lett.\ {\bf 117}, 201101 (2016); F.~Lelli, S.~S.~McGaugh, and J.~M.~Schombert,
  Astron.\ J.\ {\bf 152}, 157 (2016) (SPARC).
\bibitem{Gnedin2008} O.~Y.~Gnedin, ``Redshift evolution of the Tully--Fisher
  relation as a test of modified gravity,'' arXiv:0809.2790 (2008).
\bibitem{Skordis2021} C.~Skordis and T.~Zlosnik, Phys.\ Rev.\ Lett.\ {\bf 127},
  161302 (2021); arXiv:2007.00082.
\bibitem{Verlinde2017} E.~P.~Verlinde, SciPost Phys.\ {\bf 2}, 016 (2017).
\bibitem{Jacobson1995} T.~Jacobson, Phys.\ Rev.\ Lett.\ {\bf 75}, 1260 (1995).
\bibitem{deGraaff2024} A.~de~Graaff \textit{et al.}, JADES dynamical masses at
  $z=5.5$--$7.4$ (2024).
\bibitem{Tian2022} Y.~Tian \textit{et al.}, ``$\Lambda$CDM with baryons versus MOND:
  the time evolution of the universal acceleration scale in the Magneticum
  simulations,'' Mon.\ Not.\ R.\ Astron.\ Soc.\ {\bf 518}, 257 (2022); arXiv:2206.04333.
\bibitem{Varasteanu2025} V\u{a}r\u{a}\c{s}teanu \textit{et al.} (2025); local
  ($z<0.08$) radial-acceleration-relation acceleration scale.
\bibitem{MUSEDARK2026} MUSE--DARK collaboration, MUSE--DARK III (2026);
  $a_0$ from IFU kinematics of star-forming galaxies, $0.33<z<1.44$.
\end{thebibliography}

\end{document}
